\documentclass[11pt,a4paper]{article}
\usepackage[T1]{fontenc}
\usepackage[utf8]{inputenc}
\usepackage{amsmath,amssymb,amsthm}
\usepackage[margin=1in]{geometry}
\usepackage[colorlinks=true,linkcolor=blue,urlcolor=blue]{hyperref}
\theoremstyle{plain}
\newtheorem{theorem}{Theorem}
\newtheorem{lemma}[theorem]{Lemma}
\newtheorem{proposition}[theorem]{Proposition}
\newtheorem{corollary}[theorem]{Corollary}
\theoremstyle{definition}
\newtheorem{remark}[theorem]{Remark}

\title{Recovering the unwritten recurrences of the table OEIS A236828}
\author{Adrian Perez Fontelles\\ \small Independent researcher}
\date{2 September 2026}
\begin{document}
\maketitle

\begin{abstract}
OEIS A236828 is a two-dimensional table $T(n,k)$ whose empirical recurrences are, for several of
its lines, given only as an ORDER --- ``k=2: [order 30]'' and the like --- with the
coefficients in a linked file rather than in the entry. Those conjectures can still be settled,
and the recurrences recovered. A line of the table is a fixed-width or fixed-height array
count, hence a walk count on finitely many vertices, hence satisfies some monic recurrence of
order at most that number; Berlekamp--Massey on twice as many exact terms returns the MINIMAL
one. Where its order equals the order the entry states --- and where the same holds after the
first few terms are dropped --- any recurrence of that order the line satisfies has a
characteristic polynomial that is a multiple of the minimal one and of equal degree, hence is
that polynomial. This note settles column 2, column 3.
\end{abstract}

\noindent\small 2020 Mathematics Subject Classification. 05A15, 11B37, 15A18.\normalsize

\section{The table and the unwritten conjectures}

OEIS A236828 is ``T(n,k)=Number of (n+1)X(k+1) 0..2 arrays with the maximum plus the minimum minus the lower median of every 2X2 subblock equal.''. It is read by antidiagonals and begins
\[
81, 327, 327, 1386, 2211, 1386, 6093, 15254,\ \dots
\]
The entry carries, among its empirical claims:
\begin{quote}\small
k=2: [order 30]\\
k=3: [order 84]
\end{quote}
As of the ``Last modified'' line on the live entry (Jul 23 2025 09:36:36, revision 6) these are still
recorded as empirical, and the recurrences themselves are not in the entry text.

\section{Why a line of the table is a walk count}

Fix the index that the line fixes. The entry defines $T(n,k)$ as a number of arrays of a shape
determined by $n$ and $k$, subject to a condition local to a bounded window of consecutive
lines. Substituting the fixed index into the entry's own wording turns the definition into an
ordinary one-parameter array count, and for such a count the admissible windows are the
vertices of a finite digraph and an array is a walk. So the line is a walk count on $S$
vertices, and by the Cayley--Hamilton theorem it satisfies a monic integer recurrence of order
at most $S$.

\section{Recovering the recurrences}

\begin{lemma}\label{lem:bm}
A sequence known to satisfy some linear recurrence of order at most $S$ is determined, with its
minimal recurrence, by its first $2S$ terms; Berlekamp--Massey applied to them returns that
minimal recurrence exactly.
\end{lemma}

\subsection*{Column $k=2$}
Substituting the fixed index into the entry's wording gives an ordinary one-parameter array
count, a walk on $S=65$ vertices. Berlekamp--Massey on $2S=130$ exact terms
returns a minimal recurrence of order $30$ --- the order the entry states --- with
integer coefficients
\[
(c_1,\dots,c_{30})=(20,\ -74,\ -950,\ 7675,\ 5902,\ -208583,\ 358712,\ 2437168,\ -8245218,\ -11507689,\ 77981588,\ -13353740,\ -390127858,\ 398463912,\ 1060576100,\ -1911844548,\ -1338399672,\ 4592202520,\ -205399632,\ -6115028512,\ 2874814400,\ 4321062784,\ -3582260416,\ -1253904064,\ 1886965760,\ -92450688,\ -394737408,\ 90316800,\ 17190912,\ -4915200).
\]
so that $a(m)=\sum_{i=1}^{30}c_i\,a(m-i)$ for every $m$. Removing the first
$30+4$ terms and repeating gives order $30$ again, which is what the
identification below needs.

\subsection*{Column $k=3$}
Substituting the fixed index into the entry's wording gives an ordinary one-parameter array
count, a walk on $S=177$ vertices. Berlekamp--Massey on $2S=354$ exact terms
returns a minimal recurrence of order $84$ --- the order the entry states --- with
integer coefficients
\[
(c_1,\dots,c_{84})=(34,\ -103,\ -8609,\ 86068,\ 815732,\ -14273082,\ -24567497,\ 1251171982,\ -1875641980,\ -69522659097,\ 249574494114,\ 2625415958883,\ -14468110723020,\ -68553067413666,\ 548598961086380,\ 1178093762057447,\ -15111156938897402,\ -9582286535795279,\ 317027396438864059,\ -127453639005475181,\ -5198854880988239696,\ 6275038104321831719,\ 67616368951230545663,\ -129431066864674607728,\ -701885751998689690459,\ 1842199134597224285642,\ 5801780141253202951792,\ -20054507455432337004236,\ -37601075510334374354200,\ 173932539447048628858141,\ 182538383691400067475697,\ -1226846636654254036771934,\ -567712890836324465562644,\ 7119362338017644807732667,\ 109954859821876993999127,\ -34206718762168298395052287,\ 11871611698282821585516871,\ 136475735578962630940528860,\ -89989965590268956276346092,\ -452117035051680658878125212,\ 428378600420008851928975785,\ 1239808949588555451386976377,\ -1532064723870954677241380965,\ -2794698260754004365219828925,\ 4343288680528276372010008728,\ 5111013132962890267097372151,\ -9975030255471260423264760749,\ -7397017882009904019127916844,\ 18732818479438897085795246300,\ 8023879209490931262215825145,\ -28849302286856350040028381383,\ -5530805383617929263067292405,\ 36388937924928835098410575334,\ 221725134779149859604178876,\ -37424801669105702649176257410,\ 5402968858666800746190790160,\ 31149570381553597701315637640,\ -8422526761908424852590529656,\ -20758607923252364513095002272,\ 7889504633887572729985666816,\ 10916876213474827562358616928,\ -5245783428164741652135846976,\ -4443466584131892538570682656,\ 2575681338339155831098233792,\ 1363372746279185265494661120,\ -939143177729691102888496512,\ -303676603839791812929813248,\ 251929018675842663856442368,\ 46216933214277486436701184,\ -48808912885993272525679616,\ -4229499584030724694052864,\ 6654586152502283628118016,\ 131720903693170559156224,\ -616902042733386425237504,\ 15847028806391244947456,\ 37069230207576249729024,\ -2066803679487850708992,\ -1343393611976609103872,\ 99672213407032410112,\ 26133429260614893568,\ -2126018797931331584,\ -222129049537871872,\ 15884173650690048,\ 482990547271680).
\]
so that $a(m)=\sum_{i=1}^{84}c_i\,a(m-i)$ for every $m$. Removing the first
$84+4$ terms and repeating gives order $84$ again, which is what the
identification below needs.

\begin{theorem}
For each line above, the displayed recurrence holds for every index, and it is the recurrence
the entry records.
\end{theorem}

\begin{proof}
It holds by Lemma~\ref{lem:bm}. For the identification, the entry asserts that the line
satisfies SOME recurrence of the stated order, possibly only from some index on. Let $q$ be its
characteristic polynomial and $q_{\min}$ the minimal polynomial of the tail. Then
$q_{\min}\mid q$, and the second Berlekamp--Massey run --- on the line with its first few
terms removed --- shows $\deg q_{\min}$ equals the stated order, which is $\deg q$. Both are
monic, so $q=q_{\min}$.
\end{proof}

\section{Verification}

Three checks.

First, each line's model was compared against the entry's own published values for that line,
taken from the antidiagonal DATA read in either orientation and from the ``Table starts''
block, and agrees at every available term. This is the only point at which reading the entry's
English enters, and a misreading gives different counts.

Second, each recovered recurrence was evaluated directly on the model's values, with no
Berlekamp--Massey involved, and holds at every index tested.

Third, each order was computed twice by independent routes --- modulo a $61$-bit prime, where
every intermediate is a machine word, and again in exact rational arithmetic --- and once more
on the line with its first few terms dropped, which is what the identification requires. All
agree.

\begin{thebibliography}{9}
\bibitem{oeis} The OEIS Foundation, \emph{The On-Line Encyclopedia of Integer
Sequences}, \url{https://oeis.org}, 2026. Sequence A236828.
\bibitem{massey} J.~L.~Massey, Shift-register synthesis and BCH decoding, \emph{IEEE
Trans. Inform. Theory} \textbf{15} (1969), 122--127.
\bibitem{stanley} R.~P.~Stanley, \emph{Enumerative Combinatorics}, Volume 1, 2nd ed.,
Cambridge University Press, 2011. (Transfer-matrix method, Section 4.7.)
\bibitem{horn} R.~A.~Horn and C.~R.~Johnson, \emph{Matrix Analysis}, 2nd ed., Cambridge
University Press, 2013.
\end{thebibliography}

\end{document}
